\documentclass[12pt]{article}
\usepackage[utf8]{inputenc}
\usepackage[T1]{fontenc}
\usepackage{hyperref}
\title{Yongnan Zhuang Language Resources for K-12 Curriculum Development}
\author{Compiled by Claude Sonnet 5.5 (claude-sonnet-5-5) for Andr\'as Kornai}
\date{}
\begin{document}
\maketitle

\section{Introduction}
Yongnan Zhuang (ISO 639-3 zyn, Glottocode yong1275) is a group of related Zhuang varieties in southern China. Wikipedia describes it as a dialect-bund, not a single natural group, and cites about 1.8 million speakers (2000). Joshua Project gives about 2,031,000 speakers in China and 13,000 in Vietnam. Alternate names include Southern Zhuang, Long An and Nung An. The search found no school curricula or textbooks specific to this language, so little K-12 material is documented here.

\section{Primary Educational Resources}
No primers, textbooks or school programmes specific to Yongnan Zhuang turned up in the search results. This section is therefore empty.

\section{Linguistic and Reference Materials}
\begin{itemize}
\item Wikipedia article on the Yongnan languages: \url{https://en.wikipedia.org/wiki/Yongnan_languages}
\item Glottolog language entry (ISO zyn): \url{https://glottolog.org/resource/languoid/iso/zyn}
\item Wikitongues language page: \url{https://wikitongues.org/languages/zyn/}
\end{itemize}

\section{Religious and Community Resources}
Joshua Project reports Bible portions available as of 2023, with no complete Bible or New Testament reported.
\begin{itemize}
\item Joshua Project language page: \url{https://joshuaproject.net/index.php/languages/zyn}
\end{itemize}

\section{Digital and Audio Resources}
Global Recordings Network lists audio evangelism materials in several varieties.
\begin{itemize}
\item Global Recordings Network page, ``Zhuang, Yongnan'': \url{https://globalrecordings.net/en/language/zyn}
\item Zhuang, Southern: \url{https://www.globalrecordings.net/en/language/4979}
\item Zhuang: Qinzhou Guitai: \url{https://globalrecordings.net/en/language/6134}
\end{itemize}

\section{Curriculum Development Context}
Wikitongues rates the language EGIDS 6a (vigorous), and Glottolog lists it as not endangered. Because the group is not a single natural language, any curriculum would first have to choose which variety to base it on. The search found no existing teaching materials.

\end{document}